% GENERATED FILE. Edit canonical lessons, then run npm run generate:books.
% canonical-source-hash: fnv1a64:ff6084f8c6402747
% canonical-lessons: JA-C190-tatemono, JA-C190-shoutengai, JA-C190-hyakkaten, JA-C190-hairu, JA-C190-aku

\chapter{Buildings and Shopping Streets}
\label{ch:ja-buildings-and-shopping-streets}
\hlchaptermodality{190}

\begin{chapteropening}
\textbf{By the end of this chapter:} \emph{I can name buildings in town and say when shops open.}

Complete the last lesson of chapter 190: I can name buildings in town and say when shops open.
\end{chapteropening}

\section[\texorpdfstring{\ja{たてもの} (tatemono)}{たてもの (tatemono)}]{\ja{たてもの} (tatemono) — building}

\label{lesson:JA-C190-tatemono}



\begin{quote}
Before the new one: say the Japanese for to hand over, then the Japanese for souvenir.
\end{quote}

\subsection*{You'll want to know: \ja{たてもの}}
\textbf{\ja{たてもの}} — \emph{tatemono} — \textquotedblleft{}building\textquotedblright{}.

\begin{quote}
\textbf{\ja{あの} \ja{たかい} \ja{たてものは} \ja{ぎんこうです。}} — \emph{ano takai tatemono wa ginkō desu} — \textquotedblleft{}That tall building is a bank.\textquotedblright{}
\end{quote}

\begin{cousinweb}
\textbf{\ja{たてもの}} (\emph{tatemono}) is native Japanese, built from \emph{tate}, the stem of \emph{tateru}, \textquotedblleft{}to build\textquotedblright{}, and \emph{mono}, \textquotedblleft{}thing\textquotedblright{} (JMdict, KANJIDIC2).
\end{cousinweb}

\begin{grammarlens}[title={What you've built}]
One more word for everyday life at A2.
\end{grammarlens}

\subsection*{Your turn}
\begin{itemize}
\raggedright
  \item \emph{Say it:} \emph{tatemono}
  \item \emph{Say it:} \emph{tatemono}, once more
  \item \emph{Say it:} \emph{ano takai tatemono wa ginkō desu} — That tall building is a bank.
  \item \emph{Recall:} say the Japanese for to hand over, then the Japanese for souvenir, then say \emph{tatemono} again
\end{itemize}

\begin{tcolorbox}[breakable,colback=teal!4,colframe=teal!35!black,title={Before you move on}]
What does \ja{たてもの} mean? (\textquotedblleft{}Building\textquotedblright{}.) Say it once more.
\end{tcolorbox}
\section[\texorpdfstring{\ja{しょうてんがい} (shōtengai)}{しょうてんがい (shōtengai)}]{\ja{しょうてんがい} (shōtengai) — shopping street; shopping district}

\label{lesson:JA-C190-shoutengai}



\begin{quote}
Before the new one: say the Japanese for souvenir, then the Japanese for building.
\end{quote}

\subsection*{You'll want to know: \ja{しょうてんがい}}
\textbf{\ja{しょうてんがい}} — \emph{shōtengai} — \textquotedblleft{}shopping street; shopping district\textquotedblright{}.

\begin{quote}
\textbf{\ja{えきの} \ja{まえに} \ja{しょうてんがいが} \ja{あります。}} — \emph{eki no mae ni shōtengai ga arimasu} — \textquotedblleft{}There is a shopping street in front of the station.\textquotedblright{}
\end{quote}

\begin{cousinweb}
\textbf{\ja{しょうてんがい}} (\emph{shōtengai}) is Sino-Japanese (kango): three Chinese characters read with their on'yomi, readings borrowed from Chinese. They are \emph{shō}, \textquotedblleft{}trade\textquotedblright{} (Middle Chinese \emph{syang}), \emph{ten}, \textquotedblleft{}shop\textquotedblright{} (\emph{temH}) and \emph{gai}, \textquotedblleft{}street\textquotedblright{} (\emph{kea}) (KANJIDIC2; Middle Chinese in Baxter's transcription).
\end{cousinweb}

\begin{grammarlens}[title={What you've built}]
One more word for everyday life at A2.
\end{grammarlens}

\subsection*{Your turn}
\begin{itemize}
\raggedright
  \item \emph{Say it:} \emph{shōtengai}
  \item \emph{Say it:} \emph{shōtengai}, once more
  \item \emph{Say it:} \emph{eki no mae ni shōtengai ga arimasu} — There is a shopping street in front of the station.
  \item \emph{Recall:} say the Japanese for souvenir, then the Japanese for building, then say \emph{shōtengai} again
\end{itemize}

\begin{tcolorbox}[breakable,colback=teal!4,colframe=teal!35!black,title={Before you move on}]
What does \ja{しょうてんがい} mean? (\textquotedblleft{}Shopping street; shopping district\textquotedblright{}.) Say it once more.
\end{tcolorbox}
\section[\texorpdfstring{\ja{ひゃっかてん} (hyakkaten)}{ひゃっかてん (hyakkaten)}]{\ja{ひゃっかてん} (hyakkaten) — department store}

\label{lesson:JA-C190-hyakkaten}



\begin{quote}
Before the new one: say the Japanese for building, then the Japanese for shopping street.
\end{quote}

\subsection*{You'll want to know: \ja{ひゃっかてん}}
\textbf{\ja{ひゃっかてん}} — \emph{hyakkaten} — \textquotedblleft{}department store\textquotedblright{}.

\begin{quote}
\textbf{\ja{ひゃっかてんで} \ja{かばんを} \ja{かいます。}} — \emph{hyakkaten de kaban o kaimasu} — \textquotedblleft{}I buy a bag at the department store.\textquotedblright{}
\end{quote}

\begin{cousinweb}
\textbf{\ja{ひゃっかてん}} (\emph{hyakkaten}) is Sino-Japanese (kango): three Chinese characters read with their on'yomi, readings borrowed from Chinese. They are \emph{hyak}, \textquotedblleft{}hundred\textquotedblright{} (Middle Chinese \emph{paek}), \emph{ka}, \textquotedblleft{}goods\textquotedblright{} (\emph{xwaH}) and \emph{ten}, \textquotedblleft{}shop\textquotedblright{} (\emph{temH}) (KANJIDIC2; Middle Chinese in Baxter's transcription).
\end{cousinweb}

\begin{grammarlens}[title={What you've built}]
One more word for everyday life at A2.
\end{grammarlens}

\subsection*{Your turn}
\begin{itemize}
\raggedright
  \item \emph{Say it:} \emph{hyakkaten}
  \item \emph{Say it:} \emph{hyakkaten}, once more
  \item \emph{Say it:} \emph{hyakkaten de kaban o kaimasu} — I buy a bag at the department store.
  \item \emph{Recall:} say the Japanese for building, then the Japanese for shopping street, then say \emph{hyakkaten} again
\end{itemize}

\begin{tcolorbox}[breakable,colback=teal!4,colframe=teal!35!black,title={Before you move on}]
What does \ja{ひゃっかてん} mean? (\textquotedblleft{}Department store\textquotedblright{}.) Say it once more.
\end{tcolorbox}
\section[\texorpdfstring{\ja{はいる} (hairu)}{はいる (hairu)}]{\ja{はいる} (hairu) — to enter; to go in}

\label{lesson:JA-C190-hairu}



\begin{quote}
Before the new one: say the Japanese for shopping street, then the Japanese for department store.
\end{quote}

\subsection*{You'll want to know: \ja{はいる}}
\textbf{\ja{はいる}} — \emph{hairu} — \textquotedblleft{}to enter; to go in\textquotedblright{}.

\begin{quote}
\textbf{\ja{ともだちと} \ja{ひゃっかてんに} \ja{はいります。}} — \emph{tomodachi to hyakkaten ni hairimasu} — \textquotedblleft{}I go into the department store with a friend.\textquotedblright{}
\end{quote}

Polite \textbf{\ja{はいります}}, \emph{hairimasu}; a godan verb. It ends in -iru, but it is godan.

\begin{cousinweb}
\textbf{\ja{はいる}} (\emph{hairu}) is a native Japanese word (wago), inherited rather than borrowed; JMdict also gives a spelling that points to \emph{hau}, \textquotedblleft{}to crawl\textquotedblright{}, and \emph{iru}, \textquotedblleft{}to enter\textquotedblright{} (JMdict, KANJIDIC2).
\end{cousinweb}

\begin{grammarlens}[title={What you've built}]
One more word for everyday life at A2.
\end{grammarlens}

\subsection*{Your turn}
\begin{itemize}
\raggedright
  \item \emph{Say it:} \emph{hairu}
  \item \emph{Say it:} \emph{hairu}, once more
  \item \emph{Say it:} \emph{tomodachi to hyakkaten ni hairimasu} — I go into the department store with a friend.
  \item \emph{Recall:} say the Japanese for shopping street, then the Japanese for department store, then say \emph{hairu} again
\end{itemize}

\begin{tcolorbox}[breakable,colback=teal!4,colframe=teal!35!black,title={Before you move on}]
What does \ja{はいる} mean? (\textquotedblleft{}To enter; to go in\textquotedblright{}.) Say it once more.
\end{tcolorbox}
\section[\texorpdfstring{\ja{あく} (aku)}{あく (aku)}]{\ja{あく} (aku) — to be open; to open, to become vacant}

\label{lesson:JA-C190-aku}



\begin{quote}
Before the new one: say the Japanese for department store, then the Japanese for to enter.
\end{quote}

\subsection*{You'll want to know: \ja{あく}}
\textbf{\ja{あく}} — \emph{aku} — \textquotedblleft{}to be open; to open, to become vacant\textquotedblright{}.

\begin{quote}
\textbf{\ja{ぎんこうは} \ja{くじに} \ja{あきます。}} — \emph{ginkō wa kuji ni akimasu} — \textquotedblleft{}The bank opens at nine.\textquotedblright{}
\end{quote}

Polite \textbf{\ja{あきます}}, \emph{akimasu}; a godan verb.

\begin{cousinweb}
\textbf{\ja{あく}} (\emph{aku}) is a native Japanese word (wago), inherited rather than borrowed from Chinese; its kanji was attached later as a kun'yomi (JMdict, KANJIDIC2).
\end{cousinweb}

\begin{grammarlens}[title={What you've built}]
One more word for everyday life at A2.
\end{grammarlens}

\subsection*{Your turn}
\begin{itemize}
\raggedright
  \item \emph{Say it:} \emph{aku}
  \item \emph{Say it:} \emph{aku}, once more
  \item \emph{Say it:} \emph{ginkō wa kuji ni akimasu} — The bank opens at nine.
  \item \emph{Recall:} say the Japanese for department store, then the Japanese for to enter, then say \emph{aku} again
\end{itemize}

\begin{tcolorbox}[breakable,colback=teal!4,colframe=teal!35!black,title={Before you move on}]
What does \ja{あく} mean? (\textquotedblleft{}To be open; to open, to become vacant\textquotedblright{}.) Say it once more.
\end{tcolorbox}
